\documentclass[10pt,a4paper]{amsart}
\usepackage[margin=19mm]{geometry}
\usepackage{amsmath,amssymb,mathtools}
\usepackage[scr=rsfs,cal=euler]{mathalfa}
\usepackage{xfrac}
\usepackage[unicode,hidelinks,bookmarks=false]{hyperref}
\newcommand{\ie}{i.e.\ }
\newcommand{\cf}{cf.\ }
\newcommand{\resp}{respectively\ }
\newcommand{\Subsection}{\subsection}
\providecommand{\nobreakdash}{\nobreak\hbox{-}}
\setlength{\emergencystretch}{2em}
\multlinegap0pt
\usepackage{bm}
\usepackage{bbm}
\usepackage{dsfont}
\usepackage{xspace}
\usepackage{cancel}
\usepackage{mathtools}
\usepackage{amsthm}
\makeatletter
\@ifundefined{thm}{\newtheorem{thm}{Theorem}}{}
\@ifundefined{cor}{\newtheorem{cor}[thm]{Corollary}}{}
\@ifundefined{prop}{\newtheorem{prop}[thm]{Proposition}}{}
\makeatother
\newcommand{\R}{\mathbb{R}}
\newcommand{\rd}{\mathrm{d}}
\title[Analysis of a Biofilm Model in a Continuously Stirred Tank R]{Analysis of a Biofilm Model in a Continuously Stirred Tank Reactor with Wall Attachment}
\author{Katerina Nik, Christoph Walker}
\date{Source: arXiv:2603.10531v1 (CC BY 4.0)}
\begin{document}
\maketitle

\noindent\emph{Source and licence.} Analysis of a Biofilm Model in a Continuously Stirred Tank Reactor with Wall Attachment. Version arXiv:2603.10531v1 (CC BY 4.0), distributed under the Creative Commons Attribution 4.0 International licence, as recorded in the arXiv licence metadata for this exact version and shown on the abstract page https://arxiv.org/abs/2603.10531v1. Full rights notice: licenses/arxiv-2603.10531-CC-BY-4.0.txt.\par\smallskip

\noindent\textit{Editorial scope.} Counted unit: the Thm whose source label is \texttt{P5} in the article (source0.tex lines 920--934, source label \texttt{P5}), delivered together with the article's own complete original proof of it (source0.tex lines 936--1001, one \texttt{proof} environment of 66 lines). No mathematics is written, repaired or re-derived: statement and proof are the article's own text line for line. Required context retained uncounted: GA (lines 323--335); U (lines 351--357); P1 (lines 360--383); Cor1 (lines 501--513); H (lines 562--575); HSS (lines 872--881); HSS1 (lines 873--880); HSS2 (lines 873--880); HSS4 (lines 873--880); Q (lines 891--893); gp (lines 895--897). Every definition, display and label of those lines is retained unchanged. Omitted from this package and not redistributed: 1--322, 336--350, 358--359, 384--500, 514--561, 576--871, 881--890, 894--894, 898--919, 935--935, 1002--1593. Nothing inside a retained excerpt is omitted or altered: every retained body is delivered exactly as the article's lines are written, so no omission receipt of any kind accompanies this package. Citations: the retained excerpts carry no citation command at all, so no bibliography entry of the article is transcribed and no reference is redistributed. Reference adaptation: none was needed and none was made; every reference inside the delivered unit resolves inside it, so no unresolved reference or citation is produced. Printed slips kept literal: no printed word, symbol, hypothesis or number of the counted statement or of its proof was altered. Layout and class: the article's own class is \texttt{amsart} with packages including times, amsmath, cancel, stmaryrd, graphicx, amsfonts, enumitem, amssymb, color, mathrsfs, hyperref, cleveref, enumerate, lmodern; the wrapper is the lane's pinned \texttt{amsart} editorial wrapper at 10pt on A4, which restates the theorem-like environments the retained text needs and adds the article's own macro definitions that the retained text needs (\texttt{\textbackslash R}, \texttt{\textbackslash rd}), with extra wrapper packages (bm, bbm, dsfont, xspace, cancel, mathtools). The delivered numbering is the unit's own. Assets: no retained span contains a figure environment, a graphics-inclusion command, a PDF-page inclusion, a table or a TikZ picture; no image file is copied into this package, the package declares no asset dependency and no image file is redistributed with it. Licence: version arXiv:2603.10531v1 of this article is distributed under the Creative Commons Attribution 4.0 International licence, as recorded in the arXiv licence metadata for this exact version and shown on the abstract page https://arxiv.org/abs/2603.10531v1. A full rights notice is delivered at licenses/arxiv-2603.10531-CC-BY-4.0.txt. Characters: every retained excerpt is pure ASCII, so the delivered engine, XeLaTeX with its default Latin Modern text font, reports no missing character.
\begin{subequations}\label{GA}
\begin{equation}\label{G0}
g\in {\rm C}^{1-}(\R)\, , 
\end{equation}
where ${\rm C}^{1-}$ refers to locally Lipschitz continuity, that
\begin{equation}\label{G1}r\in
{\rm C}^{1}(\R)\ \text{  is nondecreasing with } sr(s)> 0 \text{ for }\ s\not= 0\,,
\end{equation}
and that
\begin{equation}\label{G2}
\nu,d\in {\rm C}^{1-}(\R) \ \text{ with }\ \nu(s), d(s)>  0\ \text{ for }\ s>0  \ \text{ and }\ \nu(0)=0\,.
\end{equation}
\end{subequations}

\begin{subequations}
\label{U}
\begin{align}
\kappa \partial_{y}^2u&= h^2r(u)\,, \quad 0<y<1\,,\quad \label{BBB2}\\
\partial_{y}u(0)&=0\,,\quad u(1)=S\,, \label{BBB4}
\end{align}
\end{subequations}

\begin{prop}\label{P1}
Assume \eqref{GA}. Given $h,S\in\R$, there exists a
unique solution $$u=u[h,S]\in {\rm C}^3 \bigl( [0,1]\bigr)$$ of the
boundary value problem~\eqref{U}.
It holds that
\begin{align}\label{u}
  u(y)=S-\frac{h^2}{\kappa}
  \int_y^1\int_0^\eta r \bigl( u(\xi)\bigr)\,\rd\xi \,\rd \eta,\quad y\in [0,1]\,,
\end{align}
and
$$
\Bigl[ (h,S)\mapsto u[h,S]\Bigr]\in {\rm C}^1\big(\R^2,
{\rm C}^2([0,1])\big)\, .
$$
Moreover, if $S>0$, then 
\begin{subequations}\label{ccc}
\begin{align}
&0< u(0)\le u(y)\le S\,,\quad  y\in [0,1]\,,\label{ccc1}\\
&0\le \partial_y u(y)\le \frac{h^2}{\kappa}r(u(y))\le \frac{h^2}{\kappa}r(S)\,,\quad 
y\in [0,1]\,,\label{ccc2}\\
&0\le \partial_{y}^2u(y)\le \frac{h^2}{\kappa}r(S)\,,\quad  y\in [0,1]\,. \label{ccc3}
\end{align}
\end{subequations}
\end{prop}

\begin{cor}\label{Cor1}
Assume \eqref{GA} and let $(h,S)\in [0,\infty)^2$.

{\bf (a)} The function $y\mapsto \partial_S u[h,S](y)$  is nondecreasing with respect to  $y\in [0,1]$ and
\begin{equation} \label{w1}
0<\partial_S u[h,S](0)\le \partial_S u[h,S](y)\le 1\,,\quad y\in [0,1]\,.
\end{equation}
Moreover, it holds that
$$
\partial_S\partial_y u[h,S](y)=\partial_y\partial_S u[h,S](y)\ge 0\,,\quad y\in [0,1]\,.
$$
{\bf (b)} It holds that $\partial_h u[h,S](y)\le 0$ for $y\in[0,1]$. In fact, $\partial_h u[0,S](y)=0$ for $y\in[0,1]$ and in particular $\partial_h\partial_y u[0,S](y)=0$ for $y\in[0,1]$.
\end{cor}

\begin{subequations}\label{H}
\begin{align}
\kappa \partial_{y}^2u&= h(t)^2 r(u)\,, \qquad 0<y<1\,,  && t>0\,,\label{H1}\\
\partial_{y}u(t,0)&=0,\quad u(t,1)=S(t)\,, &&t>0\,,\label{H2}\\
h'(t)&=h(t)\int_0^{1} g\big(r(u(t,y))\big)\,\rd y +\frac{\alpha}{\beta} Q(t)-d(h(t))h(t)\,,&& t>0\,, \label{H3}\\
S'(t)&=D\big(S^*-S(t)\big)-k_1 Q(t)\nu\big(S(t)\big)-\frac{\rho h(t)}{\kappa} \int_0^1  r\bigl(
u(t,y)\bigr)\,\rd y\,, && t>0\,,\label{H4}\\
Q'(t)&= \big(\nu(S(t))-k_2\big)Q(t)+\beta d\big(h(t)\big)h(t)-\alpha Q(t)\,, && t>0\,,\label{H5}
\end{align}
subject to the initial conditions
\begin{align}
h(0)&=h_0\,,\quad S(0)=S_0\,,\quad Q(0)=Q_0 \label{c6}\,.
\end{align}
\end{subequations}

\begin{subequations}\label{HSS}
\begin{align}
\kappa \partial_{y}^2u&= h^2 r(u)\,, \qquad 0<y<1\,,  \label{HSS1}\\
\partial_{y}u(0)&=0,\quad u(1)=S \,,\label{HSS2}\\
0&=h\int_0^{1} g\big(r(u(y))\big)\,\rd y +\frac{\alpha}{\beta} Q -d(h)h\,, \label{HSS3}\\
0&=D(S^*-S)-k_1 Q\nu(S)-\frac{\rho h}{\kappa} \int_0^1  r\bigl(
u(y)\bigr)\,\rd y\,, \label{HSS4}\\
0&= \big(\nu(S)-k_2\big)Q+\beta d(h)h-\alpha Q\,.\label{HSS5}
\end{align}
\end{subequations}

\begin{equation}\label{Q}
Q=\frac{\beta d(h)h}{\alpha+k_2-\nu(S)}
\end{equation}

\begin{equation}\label{gp}
\int_0^{1} g\big(r(u(y))\big)\,\rd y=\frac{\big(k_2-\nu(S)\big) d(h)}{\alpha+k_2-\nu(S)} \,.
\end{equation}

\begin{thm}\label{P5}
Assume~\eqref{GA} with  $g(0)<0$. Let $\nu\in {\rm C}^1([0,S^*])$ be increasing, and let $d\in {\rm C}^1(\R^+)$ with $\lim_{h\to\infty} d(h)h=\infty$. Moreover, assume that
\begin{equation}\label{n1}
\alpha+k_2> \nu(S^*)
\end{equation}
and that
\begin{align}\label{n2}
g(r(S^*))\big(\alpha+k_2-\nu(S^*)\big)> d(0)\big(k_2 -\nu(S^*)\big)\,.
\end{align}
Then, there is at least one nontrivial equilibrium 
$$
(h_\star,S_\star,Q_\star)\in (0,\infty)\times (0,S^*)\times (0,\infty)\,,\qquad  u_\star=u[h_\star,S_\star]\in {\rm C}^2([0,1])\big)
$$ 
to \eqref{H}. 
\end{thm}

\begin{proof} Let us first note that \eqref{n1} and the monotonicity of $\nu$ imply that
\begin{equation}\label{n3}
\alpha+k_2- \nu(S)>0\,,\quad S\in [0,S^*]\,,
\end{equation}
and recall from Proposition~\ref{P1} that
$$
\Bigl[ (h,S)\mapsto u[h,S]\Bigr]\in {\rm C}^1\big(\R^2,{\rm C}^2([0,1])\big)\,,
$$
where $u[h,S]$ is the solution to~\eqref{HSS1}-\eqref{HSS2}.
Next, owing to~\eqref{Q}
we can eliminate $Q$ and obtain from \eqref{gp} the identity
\begin{subequations}\label{Aw}
\begin{equation}\label{A1w}
G(h,S):=\int_0^{1} g\big(r(u[h,S](y))\big)\,\rd y-d(h)\left(1-\frac{\alpha}{\alpha+k_2-\nu(S)}\right)\stackrel{!}{=}0
\end{equation}
while \eqref{HSS4} yields 
\begin{equation}\label{A2}
E(h,S):=DS+\frac{\rho h}{\kappa} \int_0^1  r\bigl(u[h,S](y)\bigr)\,\rd y + \frac{k_1\beta \nu(S) d(h)h}{\alpha+k_2-\nu(S)}\stackrel{!}{=}DS^*\,.
\end{equation}
\end{subequations}
That is, \eqref{HSS} is equivalent to solving~\eqref{Aw} for $(h,S)$ and then defining $Q$ by~\eqref{Q}.

In order to solve \eqref{Aw}, let $h\ge 0$ first be fixed. Using $u[h,0]\equiv 0$, $\nu(0)=0$, $r(0)=0$, and positivity, it readily follows that
\begin{equation}\label{n6}
E(h,0)=0\,,\qquad E(h,S^*)\ge  DS^*\,.
\end{equation}
Since 
\begin{align*}
\partial_SE(h,S)&=D+\frac{\rho h}{\kappa} \int_0^1  r'\bigl(u[h,S](y)\bigr) \partial_Su[h,S](y)\,\rd y+\frac{k_1\beta \nu'(S) d(h)h(\alpha+k_2)}{(\alpha+k_2-\nu(S))^2}\,,
\end{align*}
it follows from Corollary~\ref{Cor1}~{\bf (a)}, and the monotonicity of $r$ and $\nu$ that $E(h,S)$ is strictly increasing with respect to $S$. Therefore, using~\eqref{n6} we find for each $h\ge 0$  a unique $S(h)\in (0,S^*]$ such that
\begin{equation}\label{E}
E\big(h,S(h)\big)=DS^*\,,\quad h\ge 0\,,
\end{equation}
with $S(0)=S^*$. In fact, the Implicit Function Theorem implies that 
$S\in {\rm C}^1(\R^+)$. Thus, setting $\mathsf{G}(h):=G(h,S(h))$ for $h\ge 0$, we obtain $\mathsf{G}\in {\rm C}^1(\R^+)$. Since $S(0)=S^*$ and $u[0,S^*]=S^*$, it follows from assumption~\eqref{n2} that
\begin{equation}\label{n7}
\mathsf{G}(0)=  g\big(r(S^*)\big)-d(0)\left(1-\frac{\alpha}{\alpha+k_2-\nu(S^*)}\right)>0\,.
\end{equation}
Next, the positivity of $r$ and $\nu$ yield that
$$
DS^*=E(h,S(h))\ge  \frac{k_1\beta \nu(S(h)) d(h)h}{\alpha+k_2}
$$
and since $\lim_{h\to\infty} d(h)h=\infty$, we deduce that 
\begin{equation}\label{t1}
\lim_{h\to\infty} S(h)=0\,.
\end{equation} 
Since 
$$
0\le u[h,S(h)](y)\le S(h) \rightarrow 0 
$$
for $y\in [0,1]$ as  $h\rightarrow \infty$
according to Proposition~\ref{P1}, we  conclude that
\begin{equation}\label{n4}
\lim_{h\to\infty} u[h,S(h)](y)=0\ \text{ uniformly with respect to $y\in [0,1]$}\,.
\end{equation}
 Therefore, it follows from~\eqref{n4}, \eqref{GA}, and the assumption $g(0)<0$ that there is $h_0>0$ such that $g(r(u[h_0,S(h_0)](y)))<0$ for $y\in [0,1]$. Moreover, due to~\eqref{t1} and $\nu(0)=0$ we may assume that
$$
1-\frac{\alpha}{\alpha+k_2-\nu(S(h_0))}>0\,.
$$
Hence, we obtain that
\begin{equation}\label{n8}
\mathsf{G}(h_0)\le \int_0^{1} g\big(r(u[h_0,S(h_0)](y))\big)\,\rd y <0\,.
\end{equation}
It thus follows from \eqref{n7}, \eqref{n8}, and the Mean Value Theorem that there exists some $h_\star\in (0,h_0)$ such that $\mathsf{G}(h_\star)=0$. Setting $S_\star:=S(h_\star)\in (0,S^*)$ and defining $Q_\star>0$  and $u_\star\in {\rm C}^2([0,1])$ by~\eqref{Q} and Proposition~\ref{P1}, respectively, we have established the existence of a nontrivial equilibrium for \eqref{H}.
\end{proof}

\end{document}
